\chapter{Módulo 1 --- Primeros pasos con R y RStudio}

\begin{itemize}
\tightlist
\item
  \textbf{Duración estimada:} 1.5 h teoría + 2.5 h práctica
\item
  \textbf{Objetivos de aprendizaje.} Al finalizar, la persona podrá:

  \begin{enumerate}
  \def\labelenumi{\arabic{enumi}.}
  \tightlist
  \item
    \textbf{Identificar} los paneles de RStudio (consola, editor,
    entorno, archivos/gráficos/ayuda).
  \item
    \textbf{Ejecutar} instrucciones en la consola y desde un script
    \passthrough{\lstinline!.R!}.
  \item
    \textbf{Crear} objetos con \passthrough{\lstinline!<-!} y
    \textbf{reconocer} sus tipos básicos
    (\passthrough{\lstinline!numeric!},
    \passthrough{\lstinline!integer!},
    \passthrough{\lstinline!character!},
    \passthrough{\lstinline!logical!}).
  \item
    \textbf{Usar} la ayuda (\passthrough{\lstinline!?!},
    \passthrough{\lstinline!help()!},
    \passthrough{\lstinline!example()!}) e \textbf{instalar} /
    \textbf{cargar} paquetes.
  \end{enumerate}
\end{itemize}

\section{Contenidos}

\begin{enumerate}
\def\labelenumi{\arabic{enumi}.}
\tightlist
\item
  Qué es R y qué es RStudio (el lenguaje vs.~el entorno de desarrollo).
\item
  Instalación: R desde CRAN y RStudio Desktop (Posit).
\item
  Proyectos de RStudio (\passthrough{\lstinline!.Rproj!}) y el
  directorio de trabajo.
\item
  La consola como calculadora; operadores aritméticos y de comparación.
\item
  Objetos y asignación; reglas de nombres
  (\passthrough{\lstinline!snake\_case!}).
\item
  Tipos de datos atómicos y \passthrough{\lstinline!class()!},
  \passthrough{\lstinline!typeof()!}.
\item
  Funciones: argumentos por posición y por nombre.
\item
  Paquetes: \passthrough{\lstinline!install.packages()!}
  vs.~\passthrough{\lstinline!library()!}.
\item
  Comentarios con \passthrough{\lstinline!\#!} y buenas prácticas de
  scripts.
\end{enumerate}

\section{Desarrollo}

\subsection{R como calculadora}

\begin{lstlisting}[language=R]
2 + 3 * 4        # precedencia: 14
(2 + 3) * 4      # 20
10 / 3           # 3.333333
10 %/% 3         # división entera: 3
10 %% 3          # residuo (módulo): 1
2^10             # potencia: 1024
sqrt(144)        # 12
\end{lstlisting}

\subsection{Objetos y asignación}

En R se asigna con \passthrough{\lstinline!<-!} (atajo en RStudio:
\passthrough{\lstinline!Alt!} + \passthrough{\lstinline!-!}). El nombre
va a la izquierda.

\begin{lstlisting}[language=R]
precio   <- 4999
unidades <- 12
ingreso  <- precio * unidades
ingreso            # 59988
\end{lstlisting}

\subsection{Tipos básicos}

\begin{lstlisting}[language=R]
class(3.14)        # "numeric"
class(5L)          # "integer" (la L fuerza entero)
class("Norte")     # "character"
class(TRUE)        # "logical"
is.numeric(precio) # TRUE
as.numeric("42")   # convierte texto a número: 42
\end{lstlisting}

\subsection{Funciones y argumentos}

\begin{lstlisting}[language=R]
round(3.14159, digits = 2)   # por nombre: 3.14
round(3.14159, 2)            # por posición: 3.14
seq(from = 1, to = 10, by = 3)
\end{lstlisting}

\subsection{Ayuda}

\begin{lstlisting}[language=R]
?mean              # abre la ayuda de mean()
help("seq")
example(mean)      # ejecuta los ejemplos de la ayuda
\end{lstlisting}

\subsection{Paquetes}

Un paquete se \textbf{instala una sola vez} por computadora y se
\textbf{carga en cada sesión}:

\begin{lstlisting}[language=R]
# install.packages("tidyverse")   # una vez (descomenta para instalar)
library(tidyverse)                # en cada sesión
\end{lstlisting}

\begin{quote}
\textbf{Buena práctica:} trabaja siempre dentro de un \textbf{Proyecto
de RStudio} (\passthrough{\lstinline!File → New Project!}). Así las
rutas son relativas (\passthrough{\lstinline!"datos/ventas.csv"!}) y
evitas \passthrough{\lstinline!setwd()!} con rutas absolutas que solo
funcionan en tu computadora.
\end{quote}

\section{Actividades y ejercicios}

\begin{enumerate}
\def\labelenumi{\arabic{enumi}.}
\tightlist
\item
  Crea un proyecto de RStudio llamado \passthrough{\lstinline!curso-r!}
  y un script \passthrough{\lstinline!modulo-01.R!}.
\item
  Calcula el ingreso de vender 35 audífonos de \$899 con 10 \% de
  descuento.
\item
  Guarda tu nombre en un objeto \passthrough{\lstinline!nombre!} y
  verifica su clase.
\item
  ¿Qué devuelve \passthrough{\lstinline!class(5 > 3)!}? ¿Y
  \passthrough{\lstinline!as.numeric(TRUE)!}?
\item
  Usa la ayuda para averiguar qué hace el argumento
  \passthrough{\lstinline!na.rm!} de \passthrough{\lstinline!mean()!}.
\end{enumerate}

\subsection{Soluciones}

\begin{lstlisting}[language=R]
# 2
ingreso_desc <- 35 * 899 * (1 - 0.10)
ingreso_desc                  # 28318.5

# 3
nombre <- "Ana"
class(nombre)                 # "character"

# 4
class(5 > 3)                  # "logical"
as.numeric(TRUE)              # 1

# 5: na.rm = TRUE elimina los NA antes de calcular
mean(c(10, NA, 20), na.rm = TRUE)   # 15
\end{lstlisting}

\section{Evaluación (formativa)}

Quiz de 5 preguntas:

\begin{enumerate}
\def\labelenumi{\arabic{enumi}.}
\tightlist
\item
  ¿Cuál es la diferencia entre R y RStudio?
\item
  ¿Qué operador asigna valores a un objeto en el estilo recomendado?
\item
  ¿Qué devuelve \passthrough{\lstinline!10 \%\% 4!}?
\item
  ¿Por qué \passthrough{\lstinline!install.packages()!} no se escribe en
  cada script?
\item
  Escribe una línea que calcule la raíz cuadrada de 2 redondeada a 3
  decimales.
\end{enumerate}

\emph{Respuestas:} 1) R es el lenguaje/intérprete; RStudio es el IDE. 2)
\passthrough{\lstinline!<-!}. 3) \passthrough{\lstinline!2!}. 4) Porque
solo se instala una vez; en el script basta
\passthrough{\lstinline!library()!}. 5)
\passthrough{\lstinline!round(sqrt(2), 3)!}.

\section{Referencias verificadas}

\begin{itemize}
\tightlist
\item
  \textbf{Web:} R Core Team. \emph{An Introduction to R}. CRAN.
  \url{https://cran.r-project.org/doc/manuals/r-release/R-intro.html}
  {[}EN{]} {[}OA{]}
\item
  \textbf{Libro:} Grolemund, G. (2014). \emph{Hands-On Programming with
  R}. O'Reilly Media. ISBN 978-1-4493-5901-0. Versión en línea:
  \url{https://rstudio-education.github.io/hopr/} {[}EN{]} {[}OA en
  línea{]} --- caps. 1--2.
\item
  \textbf{Libro:} Wickham, H., Çetinkaya-Rundel, M., \& Grolemund, G.
  (2023). \emph{R for Data Science} (2.ª ed.). O'Reilly Media. ISBN
  978-1-4920-9740-2. Traducción al español:
  \url{https://davidrsch.github.io/r4ds-es/} {[}ES{]} {[}OA{]} --- cap.
  ``Flujo de trabajo: conceptos básicos''.
\item
  \textbf{Video:} freeCodeCamp.org. \emph{R Programming Tutorial --
  Learn the Basics of Statistical Computing} {[}Video{]} (instructor:
  Barton Poulson). YouTube. \url{https://youtu.be/_V8eKsto3Ug} {[}EN{]}
  --- secciones ``Installing R'', ``RStudio'', ``Packages''.
  \emph{{[}fecha de publicación POR VERIFICAR{]}}
\end{itemize}
